\section{Experiments}
\label{sec:experiments}
\subsection{Anchor Configuration}
Our anchor configuration is the BioCLIP 2.5 last four blocks tile-inference submission described in Section~\ref{sec:adaptive_inference}: $4{\times}4$ tiling at tile size 448 with no overlap, an ensemble of inference at 224 and 336 pixel resolutions, softmax mean aggregation per tile, logit adjustment at $\tau{=}0.25$, and probability threshold selection at $T{=}0.03$ with a floor of two and a ceiling of ten species per quadrat. This configuration scores \textbf{0.418265} public Macro F1 on the PlantCLEF 2026 leaderboard (i002 four-blocks 224+336 ensemble; private Macro F1 $=$ 0.40283). For comparison, the same checkpoint at single-resolution 224 pixels without logit adjustment (\texttt{last\_blocks\_infer\_softmax\_mean / softmax\_mean\_probT0.05}) scores 0.41165 public / 0.38132 private --- the 224+336+LA ensemble lifts the public score by $+0.007$ and the private score by $+0.022$.
\subsection{Unfrozen-Blocks Ablation}
\label{sec:blocks_ablation}
The number of unfrozen transformer blocks is a sharp, not a plateau, design variable. We ran an ablation on the 010 recipe (the original 1.4M image manifest, five epochs, identical optimiser and head configuration), varying only the number of unfrozen blocks. Table~\ref{tab:blocks} reports the resulting public Macro F1. The head-only baseline corresponds to experiment 006 and the all-blocks (full fine-tune) row to experiment 009.
\begin{table}[ht]
\centering
\small
\begin{tabular}{lc}
\toprule
Unfrozen blocks & Public Macro F1 \\
\midrule
0 (head only) & 0.330 \\
3 blocks & 0.374 \\
\textbf{4 blocks} & \textbf{0.38333} \\
5 blocks & 0.369 \\
All 32 blocks (full FT) & 0.208 \\
\bottomrule
\end{tabular}
\caption{Unfrozen-blocks ablation on the 010 recipe. Performance peaks sharply at four blocks; one block too few costs $0.009$ Macro F1, one block too many costs $0.014$. Full fine-tuning collapses the Tree-of-Life prior and is catastrophic.}
\label{tab:blocks}
\end{table}

\begin{figure}[!htbp]
  \centering
  \input{figures/unfreeze_sweep}
  \caption{Unfrozen-blocks sweep visualised with both metrics. Held-out validation top-5 (dashed, right axis) rises monotonically with capacity from $n{=}0$ to $n{=}5$, but Kaggle Macro F1 (solid, left axis) peaks sharply at $n{=}4$ and collapses at full fine-tune ($n{=}32$). The disagreement between the two curves is the load-bearing finding of \S\ref{sec:discussion}: validation accuracy on the single-plant distribution does not predict multi-species quadrat Macro F1.}
  \label{fig:unfreeze_sweep}
\end{figure}
\subsection{Within-Backbone Saturation}
\label{sec:within_backbone}
Before pivoting to architecture-orthogonal evidence sources, we measured whether further within-backbone diversity could break the i002 anchor by comparing the four-blocks i002 checkpoint against the eight-blocks i002 checkpoint and against epoch-5 versions of both. Table~\ref{tab:within} reports the pairwise top-one agreement and submission Jaccard on 2{,}105 test quadrats.
\begin{table}[ht]
\centering
\small
\setlength{\tabcolsep}{2pt}
\begin{tabular}{@{}p{0.48\linewidth}cc@{}}
\toprule
Pair & Top-1 & Jaccard \\
\midrule
blk\_4 vs.\ blk\_8 & 0.901 & --- \\
blk\_4 vs.\ blk\_4 (epoch 5) & 0.807 & 0.664 \\
4-way intra-BC2.5 ens.\ vs.\ anchor & --- & 0.923 \\
\bottomrule
\end{tabular}
\caption{Within-backbone diagnostic gates. Top-1 agreement on a 7{,}806-way classifier and submission Jaccard at the operating threshold.}
\label{tab:within}
\end{table}
The $0.923$ Jaccard between the four way intra-BioCLIP 2.5 ensemble and our anchor submission demonstrates that all within-backbone variants are statistically equivalent on the test set: shuffling the unfrozen depth or the training-epoch index does not move the prediction distribution. This motivates the pivot to fundamentally different evidence sources.
\subsection{Pivot 1: Triple-Backbone Classifier Retraining (cRT)}
\label{sec:pivot1}
We trained two MLP classifier heads on a frozen 3328-d concatenation of BioCLIP 2.5 ViT-B, DINOv3 ViT-L, and ConvNeXt-V2-L features (head A: standard CE; head B: class-balanced random sampling). Best validation accuracy: 63.95\%. We applied the trained heads at inference using the same $4{\times}4$ tile / softmax mean aggregation pipeline as the BioCLIP 2.5 anchor.
\textbf{Pre-submission diagnostic gate.} Before submitting, we evaluated structural orthogonality against the anchor using top-one agreement and submission Jaccard, a cheap diagnostic that does not consume leaderboard submissions. Top-one agreement between the cRT head and the BioCLIP 2.5 anchor was only $0.094$ (vs.\ $0.807$ for two BioCLIP 2.5 checkpoints of the same family), and submission Jaccard at $T{=}0.05$ was $0.078$. We view this agreement-Jaccard screen as a reusable methodological gate for working-note teams: it lets a candidate ensemble be characterised as either redundant (Jaccard near 1) or structurally divergent (Jaccard near 0) before any leaderboard budget is spent.
\textbf{Submission.} Mixing the head A / head B mean of cRT with the BioCLIP 2.5 anchor at $w{=}0.3$, $T{=}0.03$, $\tau{=}0.25$ scored \textbf{0.38195} Macro F1, $\Delta = -0.036$ versus the $0.418265$ anchor.
\textbf{Reading.} cRT is structurally orthogonal but \emph{empirically wrong} on the test distribution: a peakier posterior (mean $k{=}1.33$ standalone) means high-confidence wrong picks dominate the ensemble. The $9\%$ top one agreement was diversity going \emph{against} signal rather than complementing it.
\subsection{Pivot 2: Genus / Family Co-occurrence Reranking}
\label{sec:pivot2}
For each quadrat, we computed aggregated genus and family supports
\begin{align}
P_{\text{gen}}(g) &= \!\!\sum_{s\in g}\! p_{\text{BC}}(s), \\
P_{\text{fam}}(f) &= \!\!\sum_{s\in f}\! p_{\text{BC}}(s),
\end{align}
where $p_{\text{BC}}(s)$ is the BioCLIP 2.5 anchor posterior, and applied a ``siblings stick together'' log-prior:
\begin{multline}
\log p_{\text{new}}(s) = \log p_{\text{BC}}(s) + \alpha_g \log P_{\text{gen}}(g(s)) \\ + \alpha_f \log P_{\text{fam}}(f(s)).
\end{multline}
We swept $\alpha_g \in \{0,0.1,0.2,0.3,0.5\}$, $\alpha_f \in \{0,0.05,0.1,0.2\}$, $T \in \{0.025,0.03,0.035\}$ and submitted the \emph{isovolumetric} candidate $\alpha_g{=}0.1$, $\alpha_f{=}0.0$, $T{=}0.035$, $\tau{=}0.25$ (mean $k{=}2.89$, matched to the BioCLIP 2.5 anchor volume; submission Jaccard with the anchor $=$ 0.887). Score: \textbf{0.41246}, $\Delta = -0.006$.
\textbf{Reading.} The rerank swapped 11\% of selections to genus-supported siblings at constant volume; the score moved by only $-0.006$, comfortably inside the run-to-run noise floor. Priors derived from the existing logit structure (genus support is a reweighted sum of probabilities the model already produced) carry essentially no new information --- the same evidence repackaged.
\subsection{Pivot 3: Seasonal Phenology Prior}
\label{sec:pivot3}
The full methodology is described in \S\ref{sec:phenology} and the pipeline is illustrated in Figure~\ref{fig:phenology_pipeline}. We swept $\beta \in \{0.25, 0.5, 1.0, 1.5, 2.0\}$ at three probability thresholds. The natural Bayesian operating point ($\beta{=}1.0$, no temperature scaling, $T{=}0.03$ and $\tau{=}0.25$ matching the anchor's logit-adjustment and threshold values) flips $12.8\%$ of top one picks at mean $k{=}3.24$. Score: \textbf{0.41346}, $\Delta = -0.005$.
The DOY distribution of the test set is heavily summer skewed (Jul $=$ 517, Aug $=$ 668, Apr--Jun $=$ 607 of 2{,}105 quadrats), confirming that phenology should be most informative on summer flowering taxa. Figure~\ref{fig:test_doy_wheel} visualises this concentration as a calendar wheel: roughly $85\%$ of test quadrats fall between April and August, leaving the remaining seven months almost empty.

\begin{figure}[!htbp]
  \centering
  \resizebox{0.82\linewidth}{!}{\input{figures/test_doy_wheel}}
  \caption{Calendar wheel of the PlantCLEF~2026 test set by month of observation (2{,}105 quadrats total). Radial bar length is proportional to the count for that month. The July (517) and August (668) wedges are reported exactly in the data; the April--June total of 607 is split uniformly across the three months for the figure, as is the residual 313 across the remaining seven months. The visually obvious summer over-representation motivates the seasonal phenology pivot (\S\ref{sec:pivot3}): a date-conditioned prior could only ever help in the densely-sampled months, and conversely should not be expected to lift Macro F1 on the sparse winter tail.}
  \label{fig:test_doy_wheel}
\end{figure}
\textbf{Reading.} Three concurrent factors plausibly explain the near zero delta: \textit{(i)} GBIF observation effort is itself summer biased, diluting the discriminative season signal; \textit{(ii)} the BioCLIP 2.5 backbone trained on dated images already encodes seasonal phenotype implicitly via leaf / flower / fruit state; \textit{(iii)} $11\%$ of species fall back to a uniform pdf precisely on the long tail where the visual model is most uncertain --- so phenology cannot break ties where it would be most useful.
\subsection{Inference Hyperparameter Ablations}
\label{sec:infer_ablations}
The anchor and unfrozen-blocks (\S\ref{sec:blocks_ablation}) results commit to specific values for several inference-side knobs: logit-adjustment temperature $\tau{=}0.25$, probability threshold $T{=}0.03$, softmax-mean tile aggregation, and an adaptive probability-threshold selection rule with $k_{\min}{=}2$, $k_{\max}{=}10$. This subsection consolidates the leaderboard sweeps that those values fell out of. All numbers below are public Macro F1; each table or paragraph names the underlying checkpoint, since the inference knobs were swept on both the smaller 010 anchor ($0.38333$) and the larger i002 checkpoint ($0.40590$ single-resolution, $0.418265$ as the 224+336 ensemble).

\textbf{Logit adjustment $\tau$ and probability threshold $T$.} Table~\ref{tab:tau_T} reports the joint sweep over $\tau \in \{0.25, 0.5\}$ and $T \in \{0.03, 0.05, 0.07\}$ on the single-resolution 224-pixel inference recipe (no 336-pixel ensemble). The optimum is sharp: $\tau{=}0.25$, $T{=}0.03$ dominates every other cell by between $0.022$ and $0.083$ Macro F1, and Macro F1 falls off monotonically along both axes. Doubling $\tau$ from $0.25$ to $0.50$ at fixed $T{=}0.03$ costs $0.022$ --- evidence that the long-tail boost saturates quickly and that pushing it harder starts emitting too many rare-but-implausible species.

\begin{table}[ht]
\centering
\small
\begin{tabular}{lccc}
\toprule
$\tau$ & $T{=}0.03$ & $T{=}0.05$ & $T{=}0.07$ \\
\midrule
$0.25$ & \textbf{0.40590} & 0.36423 & 0.35069 \\
$0.50$ & 0.38440 & 0.34047 & 0.32248 \\
\bottomrule
\end{tabular}
\caption{Joint $\tau \times T$ sweep on the i002 last-four-blocks checkpoint, single-resolution 224-pixel inference. Public Macro F1. The peak at $(\tau{=}0.25, T{=}0.03)$ is the anchor value taken into the 224+336 ensemble that produces the headline $0.418265$ in Table~\ref{tab:summary}.}
\label{tab:tau_T}
\end{table}

\textbf{Selection rule.} Holding $\tau{=}0.25$ fixed, we compared four selection rules with $k_{\min}{=}1$, $k_{\max}{=}10$: \emph{probability threshold} (emit every $s$ with $p_s > T$), \emph{fixed top-$k$}, \emph{gap} (emit while consecutive species differ by less than a fraction of the top probability), and \emph{relative threshold} (emit every $s$ with $p_s$ within a fraction of the top). Adaptive probability-thresholding at $T{=}0.03$ wins decisively (Table~\ref{tab:select_rule}); fixed top-$k$ is second-best at $k{=}2$, which is consistent with the test set's empirical mean of $\sim$2--3 species per quadrat; both gap- and relative-threshold rules are well off the pace.

\begin{table}[ht]
\centering
\small
\begin{tabular}{llc}
\toprule
Rule & Value & Public Macro F1 \\
\midrule
probability threshold (\emph{anchor}) & $T{=}0.03$ & \textbf{0.40590} \\
probability threshold               & $T{=}0.05$ & 0.36423 \\
probability threshold               & $T{=}0.07$ & 0.35069 \\
\midrule
fixed top-$k$                       & $k{=}1$    & 0.28633 \\
fixed top-$k$                       & $k{=}2$    & 0.39878 \\
fixed top-$k$                       & $k{=}3$    & 0.39731 \\
fixed top-$k$                       & $k{=}4$    & 0.38042 \\
fixed top-$k$                       & $k{=}5$    & 0.35959 \\
\midrule
gap                                  & $0.40$    & 0.33199 \\
gap                                  & $0.50$    & 0.33664 \\
gap                                  & $0.60$    & 0.34996 \\
\midrule
relative threshold                   & $0.15$    & 0.31348 \\
relative threshold                   & $0.20$    & 0.31704 \\
relative threshold                   & $0.25$    & 0.31823 \\
relative threshold                   & $0.30$    & 0.33623 \\
\bottomrule
\end{tabular}
\caption{Selection-rule ablation at $\tau{=}0.25$ on the i002 last-four-blocks single-resolution checkpoint. Probability-thresholding at $T{=}0.03$ dominates; fixed top-$k$ is competitive at $k{=}2$ because the test-set mean cardinality is empirically $\sim$2--3 species per quadrat. Both gap and relative-threshold variants underperform the anchor by between $0.06$ and $0.10$ Macro F1.}
\label{tab:select_rule}
\end{table}

\textbf{Per-tile aggregation} (on the 010 anchor; see ablation notes in \texttt{archive/docs/experiments\_summary.md}). We tested per-species softmax-mean against a per-species max and a mixture $\alpha \cdot \text{softmax\_mean} + (1{-}\alpha) \cdot \text{noisy\_or}$ across $\alpha \in \{0.25, 0.50, 0.75, 1.0\}$. All variants land within $0.0005$ Macro F1 of the softmax-mean anchor at $0.38333$; the second-best variant, $\alpha{=}0.75$, scores $0.38278$. We read this as evidence that the choice of per-tile aggregation is \emph{not} the bottleneck on this benchmark, and report softmax-mean throughout the rest of the paper without further reservation.

\textbf{Resolution ensemble.} Going from a single-resolution 224-pixel inference at $T{=}0.05$ to a $224+336$ probability-averaged ensemble at $\tau{=}0.25$, $T{=}0.03$ lifts public Macro F1 from $0.41165$ to $0.418265$ ($\Delta = +0.0066$) and private Macro F1 from $0.38132$ to $0.40283$ ($\Delta = +0.022$). The private gain is more than 3$\times$ the public gain --- the 336-pixel pass appears to recover rare-species pixels that the single-resolution recipe drops, and those species are the ones the private split weights most heavily.

\textbf{Takeaway.} Among the inference-side knobs, only the logit-adjustment $\tau$ and the probability threshold $T$ are load-bearing in the sense that misadjusting either by a single grid step costs more than the run-to-run noise floor of $\sim$0.002 Macro F1. The selection rule choice is qualitatively important (probability-threshold beats fixed top-$k$ by $\sim$0.007 at the optimum), but inside the probability-threshold family the surface is gentle. The per-tile aggregation is statistically flat at this scale. These observations match the broader pattern of \S\ref{sec:discussion}: the bottleneck on this benchmark is the visual model's exposure to multi-species scenes during training, not any of the inference-side hyperparameter choices.

\subsection{Aggregate Results}
\begin{table}[ht]
\centering
\scriptsize
\setlength{\tabcolsep}{2pt}
\begin{tabular}{@{}p{0.55\linewidth}cc@{}}
\toprule
Config & Macro-F1 & $\Delta$ \\
\midrule
BC2.5 last 4 blocks, 224+336+LA, $T{=}0.03$ & \textbf{0.418265} & --- \\
\hspace{1em}+ cRT triple, $w{=}0.3$ & 0.38195 & $-0.036$ \\
\hspace{1em}+ Genus $\alpha_g{=}0.1$ & 0.41246 & $-0.006$ \\
\hspace{1em}+ Pheno., $\beta{=}1.0$ & 0.41346 & $-0.005$ \\
\bottomrule
\end{tabular}
\caption{All four submitted configurations on the PlantCLEF 2026 hidden test set. The cRT pivot is strongly negative; the genus rerank and the phenology prior move the anchor by $-0.006$ and $-0.005$ respectively --- magnitudes comfortably inside the run-to-run leaderboard noise floor we previously characterised at ${\sim}0.002$ but consistent in sign. We read this as the BioCLIP 2.5 posterior implicitly absorbing both signals through training-distribution co-occurrence: neither prior carries enough orthogonal information to break the visual anchor.}
\label{tab:summary}
\end{table}

\begin{figure}[!htbp]
  \centering
  \input{figures/aggregate_deltas}
  \caption{Public Macro F1 of the four submitted configurations from Table~\ref{tab:summary}. The dashed grey line marks the anchor (0.418265). cRT is the only pivot whose loss exceeds the characterised noise floor; the genus and phenology priors are both within $-0.006$ of the anchor.}
  \label{fig:aggregate_deltas}
\end{figure}
